%% kapa: exam=1 round=23 year=2012 period=1 form=A subject=civil_law subject_ko=민법 seq=19 number=19
%% source: Q-Net 감정평가사 기출문제 (cst00309) artlSeq=5052578 file=raw/1008/5052578/2012년 제23회 감정평가사 제1차시험 기출문제(1교시 A형).pdf
%% conversion: kordoc 4.24.51 -> tools/gen1.py (auto); review=proofread-sonnet; flags=none
\begin{question}{19}
甲은 乙을 강박하여 乙 소유의 X 토지를 무상으로 자신에게 증여하도록 강요하였다. 이에 乙은 자신의 재산을 甲에게 강제로 빼앗긴다는 인식을 하였지만 그래도 어쩔 수 없다고 생각하면서 X 토지를 무상으로 증여하였다. 다음 설명 중 옳은 것은? (다툼이 있으면 판례에 의함)
\begin{choices}
  \item 乙의 증여의 의사표시는 비진의 의사표시에 해당한다.
  \item 甲의 불법적인 해악의 고지로 말미암아 乙이 공포심을 느꼈다면 그것이 乙의 의사결정에 큰 영향을 주지 않았더라도 乙의 증여행위는 무효가 된다.
  \item 불공정한 법률행위를 이유로 乙의 증여행위를 무효로 할 수 없다.
  \item 乙은 강박에 의한 의사표시를 이유로 증여행위를 취소할 수 없다.
  \item 乙의 증여의 의사표시는 통정허위표시에 해당하여 무효이다.
\end{choices}
\end{question}
\answer{3}
